\documentclass[11pt]{article}
\usepackage[a4paper,margin=21mm]{geometry}
\usepackage[no-math]{fontspec}
\setmainfont{Latin Modern Roman}
\newfontfamily\CJKfont{Noto Serif CJK SC}
\XeTeXlinebreaklocale "zh"
\XeTeXlinebreakskip=0pt plus 1pt
\usepackage{amsmath,amssymb,amsthm,amscd,graphicx,color}
\usepackage{xurl}
\usepackage[unicode,hidelinks]{hyperref}
\allowdisplaybreaks
\setlength\emergencystretch{3em}
\numberwithin{equation}{section}

\newcommand{\RS}[1]{\textcolor{red}{#1}}

\numberwithin{equation}{section}

\newcommand{\dd}{\mathrm{d}}
\newcommand{\Res}{\mathop{\,\rm Res\,}}
\DeclareMathOperator*{\res}{Res}
\newcommand{\bs}[1]{\ensuremath{\boldsymbol{#1}}}
\newcommand{\Pf}{{\rm Pf}}
\newcommand{\Tr}{{\rm Tr}}
\renewcommand{\arraystretch}{1.2}

\usepackage{amsfonts,esint}

\newtheorem{Theorem}{Theorem}[section]
\newtheorem*{Theorem*}{Theorem}
\newtheorem{Corollary}[Theorem]{Corollary}
\newtheorem{Lemma}[Theorem]{Lemma}
\newtheorem{Proposition}[Theorem]{Proposition}
 { \theoremstyle{definition}
\newtheorem{Definition}[Theorem]{Definition}
\newtheorem{Note}[Theorem]{Note}
\newtheorem{Example}[Theorem]{Example}
\newtheorem{Remark}[Theorem]{Remark} }

\begin{document}
\begin{center}\Large Regularized Pfaffian formula for even-dimensional Kontsevich Hermitian matrix integrals with a positive external matrix and a continuous all-moments potential\end{center}
\noindent\textbf{Stable ID:} 988\quad\textbf{Author:} Gaëtan Borot; Raimar Wulkenhaar\par
\noindent\textbf{Work:} A Note on BKP for the Kontsevich Matrix Model with Arbitrary Potential\par
\noindent\textbf{Source:} \url{https://sigma-journal.com/2024/050/}\par
\noindent\textbf{Version:} Official SIGMA 20 (2024) article 050, DOI 10.3842/SIGMA.2024.050; journal PDF and native TeX acquired 2026-09-30, exact bytes pinned by SHA256\par
\noindent\textbf{Rights:} CC BY 4.0. Authors retain copyright. \url{https://creativecommons.org/licenses/by/4.0/}. Reuse and typesetting adaptation; original language and mathematical content preserved.\par
\noindent\textbf{Difficulty (prior selection preserved):} {\CJKfont H2: The central obligation is to extend the de Bruijn Pfaffian integration identity through a singular anti-diagonal kernel while controlling the noncompact integral. The author factors the determinant with Vandermonde divided differences and a simplex derivative formula to construct an integrable domination uniform in the cutoff. Tracking the matrix-integral normalization and absorbing the external Vandermonde then gives the stated exact Pfaffian formula.}\par
\medskip\noindent\textbf{Editorial boundary note (not author text):} Complete original Theorem1.1, even-N regularized Pfaffian formula, with the exact Hermitian Lebesgue and Gaussian normalization, positive external matrix, continuous all-moments potential and formal odd deformation parameters. The complete current proof retains volume/HCIZ normalization, singular de Bruijn extension lemma and full divided-difference/simplex integrable domination, cutoff limit, Pfaffian conjugation, Taylor remainders and every prefactor. Original principal-value integral notation uses the real installed esint package. The Pfaffian evaluation is a substantial explicit matrix-integral outcome; its singular-kernel lemma remains supporting core and is not counted separately. The later BKP bilinear identity and examples are excluded. Standard HCIZ, Schur Pfaffian, classical de Bruijn and simplex divided-difference identities remain precise author references; no generated verification, new proof or formula repair. Literal source conventions and printed indices are preserved.\par
\medskip\hrule\medskip\noindent\textbf{Author text begins below.}\par

% Original source lines 69--122
\section{The formula}\label{Chap1}

\subsection{The Kontsevich matrix model with arbitrary potential}

Let $\mathcal{H}_N$ be the space of Hermitian $N \times N$ matrices
equipped with the Lebesgue measure
\[
\dd H= \prod_{i=1}^N \dd H_{ii}
 \prod_{1 \leq i < j \leq N} \dd\operatorname{Re}(H_{ij})\, \dd\operatorname{Im}(H_{ij}).
\]
Given a positive matrix $\Lambda = \operatorname{diag}(\lambda_1,\dots,\lambda_N)$,
we introduce the Gaussian probability measure on $\mathcal{H}_N$
\begin{equation}
\label{theGauss}
\dd \mathbb{P}_N(H) = \frac{\sqrt{\Delta(\boldsymbol{\lambda},\boldsymbol{\lambda})}}{2^{\frac{N}{2}} (2\pi)^{\frac{N^2}{2}}}\,\dd H {\rm e}^{-\frac{1}{2}{\rm Tr}(\Lambda H^2)}.
\end{equation}
We use the notations $\Delta(\boldsymbol{\lambda}) = \prod_{1 \leq i < j \leq N} (\lambda_j - \lambda_i)$ and $\Delta(\boldsymbol{\lambda},\boldsymbol{\mu}) = \prod_{1 \leq i,j \leq N} (\lambda_i + \mu_j)$. We denote $[n] = \{1,\dots,n\}$ and $\lambda_{\min} = \min\{\lambda_i\mid i \in [N]\}$.

Let $V_0$ be a continuous function on $\mathbb{R}$ such that the measure ${\rm e}^{- \frac{1}{2}\lambda_{\min} x^2 + V_0(x)}\,\dd x$ has finite moments on $\mathbb{R}$ (take for instance $V_0$ to be a polynomial of even degree with negative top coefficient). Then the measure $\dd \mathbb{P}_N(H) {\rm e}^{ \operatorname{Tr}V_0(H)}$ on $\mathcal{H}_N$ is finite. Let
\[
V_{\mathbf{t}}(x) = V_0(x) + \sum_{k \geq 0} t_{2k + 1} x^{2k + 1},
\]
where $\mathbf{t} = (t_{2k + 1})_{k \geq 1}$ are formal parameters. The partition function of the Kontsevich model with arbitrary potential is defined by
\begin{equation}
\label{Zdef} Z_N(\mathbf{t}) = \int_{\mathcal{H}_N} \dd\mathbb{P}_N(H) {\rm e}^{\operatorname{Tr}V_{\mathbf{t}}(H)}.
\end{equation}
This note aims at showing that $Z_N(\mathbf{t})$ is a
tau-function of the BKP hierarchy \cite{Date:1981qz}.

\subsection{Pfaffian formula}

First, we establish a Pfaffian formula for $Z_N(\mathbf{t})$. The proof proposed in Section~\ref{S2} consists in classical algebraic manipulations with matrix integrals and an analysis argument -- that one may find of independent interest, see Lemma~\ref{HUHH} and Remark~\ref{therem} -- to justify the extension of de~Bruijn's Pfaffian formula \cite{deBruijn1955} to singular kernels like the one appearing in \eqref{BZ}.
\begin{Theorem}
\label{th:1} For even $N$, we have
\begin{equation}
\label{BZ} Z_N(\mathbf{t}) = \frac{\sqrt{\Delta(\boldsymbol{\lambda},\boldsymbol{\lambda})}}{2^{\frac{N^2}{2}} (2\pi)^{\frac{N}{2}}\prod_{n = 1}^{N - 1} n!} \operatorname{Pf}_{0 \leq m,n \leq N - 1}(K_{N;m,n}(\mathbf{t})),
\end{equation}
where $\fint = \lim_{\epsilon \rightarrow 0} \int_{|x + y| \geq \epsilon}$ is the Cauchy principal value integral and
\begin{gather*}
 K_{N;m,n}(\mathbf{t}) = \fint_{\mathbb{R}^2} \frac{x - y}{x + y} F_{N;m}(x)F_{N;n}(y) {\rm e}^{V_{\mathbf{t}}(x) + V_{\mathbf{t}}(y)}\,\dd x \dd y, \\
F_{N;n}(x) = x^{2n} + \frac{(-2)^n n!}{\Delta(\boldsymbol{\lambda})}\det\bigl(\lambda_i^0 \big| \lambda_i^1 \big| \cdots \big| \lambda_i^{n - 1} \big| R_N\bigl(-\tfrac{1}{2}\lambda_ix^2\bigr) \big| \lambda_i^{n + 1} \big| \cdots \big| \lambda_i^{N - 1}\bigr), \\
R_N(\xi) = \frac{\xi^N}{(N - 1)!} \int_{0}^{1} \dd u\, (1 - u)^{N - 1} {\rm e}^{\xi u} = {\rm e}^{\xi}\biggl(1 - \frac{\Gamma(N;\xi)}{(N - 1)!}\biggr).
\end{gather*}
\end{Theorem}

Note that $R_N(\xi)$ is simply the $N$-th remainder of the Taylor series of ${\rm e}^{\xi}$. The expression given for $F_{N;n}(x)$ emphasises that it is $x^{2n} + O\bigl(x^{2N + 2}\bigr)$ as $x \rightarrow 0$, but we also have the equivalent expression
\[
F_{N;n}(x) = \frac{(-2)^n n!}{\Delta(\boldsymbol{\lambda})} \det\bigl(\lambda_i^0 \big| \lambda_i^1 \big| \cdots \big| \lambda_i^{n - 1} \big| {\rm e}^{-\frac{1}{2}\lambda_i x^2} \big| \lambda_i^{n + 1} \big| \cdots \big| \lambda_i^{N - 1}\bigr).
\]
For instance, the formula for $N = 2$ involves the two functions
\[
F_{2;0}(x) = \frac{{\rm e}^{-\frac{1}{2}\lambda_1 x^2} \lambda_2 - {\rm e}^{-\frac{1}{2}\lambda_2 x^2}\lambda_1}{\lambda_2 - \lambda_1}, \qquad F_{2;1}(x) = 2 \frac{{\rm e}^{-\frac{1}{2}\lambda_1 x^2} - {\rm e}^{-\frac{1}{2}\lambda_2 x^2}}{\lambda_2 - \lambda_1} .
\]


\setcounter{section}{1}
% Original source lines 195--343
\section{Proof of Theorem~\ref{th:1}}
\label{S2}
A Hermitian matrix $H$ decomposes as
\begin{gather*}
H = \sum_{a = 1}^N (\operatorname{Re} H_{a,a}) E_{a,a}+ \sum_{1 \leq a < b \leq N} \bigl(\sqrt{2} \operatorname{Re} H_{a,b}\bigr) \tfrac{1}{\sqrt{2}}(E_{a,b} + E_{b,a}) \\
\hphantom{H =}{}
 + \bigl(\sqrt{2} \operatorname{Im} H_{a,b}\bigr) \tfrac{{\rm i}}{\sqrt{2}}(E_{a,b} - E_{b,a}).
\end{gather*}
As $E_{a,a}$, $\frac{1}{\sqrt{2}}(E_{a,b} + E_{b,a})$ and $\frac{{\rm i}}{\sqrt{2}}(E_{a,b} - {\rm E}_{b,a})$ have unit norm for the standard Euclidean metric on ${\rm Mat}_{N}(\mathbb{C}) \cong \mathbb{R}^{2N^2}$, the volume form on $\mathcal{H}_N$ induced by the Euclidean volume form on ${\rm Mat}_N(\mathbb{C}) \cong \mathbb{R}^{2N^2}$ is $2^{\frac{N(N - 1)}{2}}\,\dd H$. Denote $\mathcal{U}_N$ the unitary group and $\dd\nu$ its volume form induced by the Euclidean volume form in ${\rm Mat}_N(\mathbb{C})$. The corresponding volume is
 \[
 {\rm Vol}(\mathcal{U}_N) = \frac{(2\pi)^{\frac{N(N + 1)}{2}}}{\prod_{n = 1}^{N - 1} n!}.
\]
We also recall the Harish-Chandra--Itzykson--Zuber formula \cite{Itzykson:1979fi}
\[
\frac{1}{\prod_{n = 1}^{N - 1} n!} \int_{\mathcal{U}_N} \frac{\dd \nu(U)}{{\rm Vol}(\mathcal{U}_N)} {\rm e}^{{\rm Tr}(AUBU^{\dagger})} = \frac{\det\bigl({\rm e}^{a_ib_j}\bigr)}{\Delta(\boldsymbol{a})\Delta(\boldsymbol{b})},
\]
 where $A = \operatorname{diag}(a_1,\dots,a_N)$ and $B = \operatorname{diag}(b_1,\dots,b_N)$.

Diagonalising the matrix $H = UXU^{\dagger}$ with $X = \operatorname{diag}(x_1,\dots,x_N)$ and $U \in \mathcal{U}_N$ defined up to action of $\mathfrak{S}_N \times \mathcal{U}_1^{N}$ brings the partition function \eqref{Zdef} in the form
\begin{align}
Z_N(\mathbf{t}) & = \frac{\sqrt{\Delta(\boldsymbol{\lambda},\boldsymbol{\lambda})}}{2^{\frac{N}{2}} (2\pi)^{\frac{N^2}{2}}} \frac{1}{N! (2\pi)^N 2^{\frac{N(N - 1)}{2}}} \nonumber\\
&\quad{}\times \int_{\mathbb{R}^N} \biggl(\int_{\mathcal{U}_N} \dd \nu(U) {\rm e}^{-\frac{1}{2}{\rm Tr}(\Lambda U X^2 U^{\dagger})}\biggr) (\Delta(\boldsymbol{x}))^2 \prod_{i = 1}^N {\rm e}^{V_{\mathbf{t}}(x_i)}\dd x_i \nonumber \\
& = \frac{\sqrt{\Delta(\boldsymbol{\lambda},\boldsymbol{\lambda})}}{2^{\frac{N^2}{2}} (2\pi)^{\frac{N}{2}} N! \Delta\bigl(-\frac{\boldsymbol{\lambda}}{2}\bigr)} \int_{\mathbb{R}^N} \frac{(\Delta(\boldsymbol{x}))^2}{\Delta\bigl(\boldsymbol{x}^2\bigr)} \mathop{\det}_{1 \leq i,j \leq N} \bigl({\rm e}^{-\frac{1}{2}\lambda_ix_j^2}\bigr) \prod_{i = 1}^N {\rm e}^{V_{\mathbf{t}}(x_i)} \dd x_i. \label{Ztint}
\end{align}
Here we could use Fubini because the integrand in the first line of \eqref{Ztint} is real positive, and in fact integrable due to the assumptions on $V_0$ and $\Lambda$. We observe that \smash{$\Delta\bigl(-\frac{\boldsymbol{\lambda}}{2}\bigr) = (-2)^{-\frac{N(N - 1)}{2}}\Delta(\boldsymbol{\lambda})$} and recall Schur's Pfaffian identity \cite{Schur1911}, for $N$ even
\[
\frac{(\Delta(\boldsymbol{x}))^2}{\Delta\bigl(\boldsymbol{x}^2\bigr)} = \prod_{1 \leq i < j \leq N} \frac{x_j - x_i}{x_j + x_i} = \operatorname{Pf}_{1 \leq i,j \leq N}\biggl(\frac{x_j - x_i}{x_j + x_i}\biggr).
\]
So, up to a prefactor, $Z_N(\mathbf{t})$ is an integral of the form
\begin{equation}
\label{Sint}
\int_{\mathbb{R}^N} \operatorname{Pf}_{1 \leq i,j \leq N} (S(x_i,x_j)) \mathop{{\rm det}}_{\substack{0 \leq m \leq N - 1 \\ 1 \leq j \leq N}} \bigl(f_m\bigl(x_j^2\bigr)\bigr) \prod_{i = 1}^N \rho(x_i)\,\dd x_i.
\end{equation}
De Bruijn's identity \cite{deBruijn1955} would allow rewriting \eqref{Sint} as
\begin{equation}
\label{dJ}
N! \operatorname{Pf}_{0 \leq m,n \leq N - 1}\biggl(\int_{\mathbb{R}^2} S(x,y) f_m\bigl(x^2\bigr)f_n\bigl(y^2\bigr)\rho(x)\rho(y)\,\dd x \dd y\biggr),
\end{equation}
but the proof in loc.\ cit.\ is solely based on algebraic manipulations, valid when $(f_n)_{n = 0}^{N - 1}$ is a~sequence of measurable functions on $\mathbb{R}_{\geq 0}$ and $S(x,y) = -S(y,x)$ is a measurable function on $\mathbb{R}^2$ such that $\int_{\mathbb{R}^2} \big|S(x,y)f_m\bigl(x^2\bigr)f_n\bigl(y^2\bigr)\big|\rho(x)\rho(y)\dd x\dd y < +\infty$. The choice of $S(x,y) = \frac{x - y}{x + y}$ in general violates this integrability assumption due to the presence of the simple pole on the anti-diagonal combined with the non-compactness of $\mathbb{R}^2$. Nevertheless, we show that the conclusion~\eqref{dJ} remains valid provided the integral in the Pfaffian is understood as a Cauchy principal value, under a Schwartz-type condition.

\begin{Lemma}
\label{HUHH}Let $\rho > 0$ be a measurable function on $\mathbb{R}$ and $\mathbb(f_n)_{n = 0}^{N - 1}$ be a sequence of $\mathcal{C}^{N - 1}$-functions on $\mathbb{R}_{\geq 0}$ such that $f^{(\ell)}_m$ is bounded by a polynomial for any $m,\ell \in \{0,\dots,N - 1\}$. Let $S(x,y) = \frac{\tilde{S}(x,y)}{x + y}$ where $\tilde{S}$ is a measurable function on $\mathbb{R}^2$ such that
\[
\forall k,l \in \mathbb{Z}_{\geq 0}\qquad \int_{\mathbb{R}^2} \bigl|\tilde{S}(x,y)x^ky^l\bigr|\rho(x)\rho(y)\,\dd x \dd y < +\infty.
\]
Then, for $N$ even
\begin{gather*}
 \int_{\mathbb{R}^N} \operatorname{Pf}_{1 \leq i,j \leq N} (S(x_i,x_j)) \mathop{\det}_{\substack{0 \leq m \leq N - 1 \\
 1 \leq j \leq N}} \bigl(f_m\bigl(x_j^2\bigr)\bigr) \prod_{i = 1}^n \rho(x_i)\dd x_i \\
\qquad\quad{} = N!\operatorname{Pf}_{0 \leq m,n \leq N - 1}\biggl(\fint_{\mathbb{R}^2} S(x,y) f_m\bigl(x^2\bigr) f_n\bigl(y^2\bigr)\rho(x)\rho(y)\,\dd x \dd y\biggr),
\end{gather*}
where $\fint = \lim_{\epsilon \rightarrow 0} \int_{|x + y| \geq \epsilon}$ and the integrand in the left-hand side is integrable.
\end{Lemma}
\begin{proof} Take $\epsilon > 0$ and set $S_{\epsilon}(x,y) = S(x,y) \cdot \mathbf{1}_{|x + y| \geq \epsilon}$. In this situation, we can use de Bruijn's formula and write
\begin{gather}
\nonumber \int_{\mathbb{R}^N} \operatorname{Pf}_{1 \leq i,j \leq N} (S_{\epsilon}(x_i,x_j)) \mathop{\det}_{\substack{0 \leq m \leq N - 1 \\ 1 \leq j \leq N}} \bigl(f_m\bigl(x_j^2\bigr)\bigr) \prod_{i = 1}^N \rho(x_i)\dd x_i \\
\qquad {} = N!\operatorname{Pf}_{0 \leq m,n \leq N - 1} \biggl(\int_{\mathbb{R}^2} S_{\epsilon}(x,y)f_m\bigl(x^2\bigr)f_n\bigl(y^2\bigr)\rho(x)\rho(y)\,\dd x \dd y\biggr).\label{LHs}
\end{gather}
The right-hand side tends to
\[
N! \operatorname{Pf}_{0 \leq m,n \leq N - 1} \biggl(\fint_{\mathbb{R}^2} S(x,y) f_m\bigl(x^2\bigr)f_n\bigl(y^2\bigr)\rho(x)\rho(y)\dd x \dd y\biggr)
\]
 when ${\epsilon \rightarrow 0}$. Call $I_{\epsilon}(\mathbf{x})$ the integrand in the left-hand side of formula~\eqref{LHs}. We clearly have $\lim_{\epsilon \rightarrow 0} I_{\epsilon}(\mathbf{x}) = I_0(\mathbf{x})$ for $\mathbf{x}$ almost everywhere in $\mathbb{R}^N$. Provided we can find for $I_{\epsilon}(\mathbf{x})$ a uniform in $\epsilon$ and integrable on~$\mathbb{R}^2$ upper bound, the lemma follows from dominated convergence.

To find such a bound, we introduce the matrix $W(\boldsymbol{\xi})$ with entries $\xi_j^{n}$ at row index $n \in \{0,\dots,N - 1\}$ and column index $j \in [N]$, which satisfies $\Delta(\boldsymbol{\xi}) = \det W(\boldsymbol{\xi})$. Its inverse matrix is
\[
\bigl(W(\boldsymbol{\xi})^{-1}\bigr)_{i,n} = \frac{(-1)^{N - n - 1}e_{N - n - 1}\bigl(\boldsymbol{\xi}_{[i]}\bigr)}{\prod_{j \neq i} (\xi_i - \xi_j)},
\]
where $e_k$ is the $k$-th elementary symmetric polynomial and $\boldsymbol{\xi}_{[i]} = \bigl(\xi_1,\dots,\widehat{\xi_i},\dots,\xi_N\bigr)$. Then
\begin{gather}
\nonumber \mathop{\det}_{\substack{0 \leq m \leq N - 1 \\ 1 \leq j \leq N }} \bigl(f_m\bigl(x_j^2\bigr)\bigr) = \Delta\bigl(\boldsymbol{x}^2\bigr) \det_{\substack{1 \leq i \leq N \\ 0 \leq n \leq N - 1}} \bigl(f_n\bigl(x_i^2\bigr)\bigr) \cdot \det \bigl(W\bigl(\boldsymbol{x}^2\bigr)^{-1}\bigr) \\
\hphantom{\mathop{\det}_{\substack{0 \leq m \leq N - 1 \\ 1 \leq j \leq N }} \bigl(f_m\bigl(x_j^2\bigr)\bigr)}{} = \Delta\bigl(\boldsymbol{x}^2\bigr) \det_{0 \leq m,n \leq N - 1} \Biggl(\sum_{i = 1}^{N} \frac{(-1)^{N - m - 1} f_n\bigl(x_i^2\bigr) e_{N - m - 1}\bigl(\boldsymbol{x}^2_{[i]}\bigr)}{\prod_{j \neq i} \bigl(x_i^2 - x_j^2\bigr)}\Biggr).\label{26}
\end{gather}
Up to a sign that we can take out of the determinant, the $(m,n)$-entry inside the determinant is
\begin{align*}
 \bigl[u^{N - m - 1}\bigr] \sum_{i = 1}^N f_n\bigl(x_i^2\bigr) \prod_{j \neq i} \frac{1 + ux_j^2}{x_i^2 - x_j^2} & = \bigl[u^{N - m - 1}\bigr] \prod_{i = 1}^N \bigl(1 + ux_i^2\bigr)\!\Biggl( \sum_{i = 1}^N \frac{f_n\bigl(x_i^2\bigr)}{1 + ux_i^2} \frac{1}{\prod_{j \neq i} (x_i^2 - x_j^2)}\Biggr) \\
& = \sum_{k = 0}^{N - m - 1} e_{N - m - 1 - k}\bigl(\boldsymbol{x}^2\bigr)\!\Biggl( \sum_{i = 1}^{N} \frac{(-1)^k x_i^{2k} f_{n}\bigl(x_i^2\bigr)}{\prod_{j \neq i} \bigl(x_i^2 - x_j^2\bigr)}\Biggr).
\end{align*}
In the first two steps, $[u^m]$ acting on the formal power series of $u$ to its right meant extracting the coefficient of $u^m$. Up to the use of squared variables, we recognise the divided difference
\[
g[\xi_1,\dots,\xi_N] := \sum_{i = 1}^N \frac{g(\xi_i)}{\prod_{j \neq i} (\xi_i - \xi_j)}.
\]
When $g$ is $\mathcal{C}^{N - 1}$, it can be written (see, e.g., \cite[Theorem 2, p.~250]{AnaN}) as an integral over the $(N - 1)$-dimensional simplex $\Delta_{N - 1} = \{p \in [0,1]^N \mid p_1 + \cdots + p_N = 1\}$, equipped with the volume form $\dd \sigma(\boldsymbol{p}) = \dd p_1 \cdots \dd p_{N - 1}$:
\[
g[\xi_1,\dots,\xi_N] = \int_{\Delta_{N - 1}} g^{(N - 1)}(p_1\xi_1 + \cdots + p_N\xi_N)\,\dd \sigma(\boldsymbol{p}).
\]
We use this for $g_{k,n}(\xi) = (-1)^{k}\xi^k f_n(\xi)$. Inserting the integral representation in \eqref{26} yields
\begin{gather*}
|I_{\epsilon}(\mathbf{x})| = \bigl|\Delta\bigl(\boldsymbol{x}^2\bigr)\bigr| \operatorname{Pf}_{1 \leq i,j \leq N}(S_{\epsilon}(x_i,x_j) ) \\
 \times \Biggl| \mathop{{\rm det}}_{0 \leq m,n \leq N - 1}\!\Biggl( \sum_{k = 0}^{N - 1 - m}\! e_{N - 1 - m - k}\bigl(\boldsymbol{x}^2\bigr) \int_{\Delta_{N - 1}} g_{k,n}^{(N - 1)}\bigl(p_1x_1^2 + \cdots + p_N x_N^2\bigr) \dd \sigma(\boldsymbol{p})\Biggr)\Biggr| \prod_{i = 1}^N \rho(x_i).
\end{gather*}
Since $\bigl|\Delta\bigl(\boldsymbol{x}^2\bigr)\bigr|$ cancels the denominators in $S_{\epsilon}$, the first line of the right-hand side admits an upper bound
by sum of terms, each of which is a polynomial in $\boldsymbol{x}$ multiplied by \smash{$\prod_{\{i,j\} \in \mathcal{P}} \bigl|\tilde{S}(x_i,x_j)\bigr|$}, where $\mathcal{P}$ is a partition of $[N]$ into pairs. In the second line, we first expand the determinant inside the absolute value and use the triangular inequality to get an upper bound by a sum of finitely many positive terms, each of which involves an $N$-fold product of simplex integrals of functions with at most polynomial growth, since the derivatives $f_{n}^{(\ell)}$ \big(and thus $g_{k,n}^{(N - 1)}$\big) have at most polynomial growth. Therefore, they result in a polynomial upper bound in the variable $\boldsymbol{x}$. We are thus left with an upper bound by a sum of finitely many terms of the form
\[
\prod_{\{i,j\} \in \mathcal{P}} \bigl|\tilde{S}(x_i,x_j)\bigr| \prod_{i = 1}^{N} x_i^{q_i}\rho(x_i) = \prod_{\{i,j\} \in \mathcal{P}} \bigl|\tilde{S}(x_i,x_j)\bigr|x_i^{q_i}x_j^{q_j} \rho(x_i)\rho(x_j)
\]
for various $N$-tuples of integers $\boldsymbol{q}$ and pair partitions $\mathcal{P}$ of $[N]$. Integrating each term of this form over $\mathbb{R}^{N}$ factorizes into a product of $\frac{N}{2}$ two-dimensional integrals, each of them being finite by assumption. This provides the domination assumption to conclude $\smash{\lim_{\epsilon \rightarrow 0} \int_{\mathbb{R}^N} I_{\epsilon}(\boldsymbol{x}) \prod_{i = 1}^N \dd x_i} = \smash{\int_{\mathbb{R}^N} I_0(\boldsymbol{x}) \prod_{i = 1}^N \dd x_i}$ as desired.
\end{proof}

\begin{Remark}
\label{therem} The proof can easily be adapted to obtain an analogous statement for kernels of~the form \smash{$S(x,y) = \frac{\tilde{S}(x,y)}{x - y}$}, in which case one can use $f_n(x)$ instead of $f_n\bigl(x^2\bigr)$.
\end{Remark}

The assumptions of Lemma~\ref{HUHH} are fulfilled for
\[
S(x,y) = \frac{x - y}{x + y},\qquad \rho(x) = {\rm e}^{-\frac{1}{2}\lambda_{\min} x^2 + V_{\mathbf{t}}(x)},\qquad f_m(\xi) = {\rm e}^{-\frac{1}{2}(\lambda_{m+1} - \lambda_{\min})\xi},
\]
where we stress that $\mathbf{t}$ are formal parameters. Therefore, coming back to \eqref{Ztint} and tracking the $N$-dependent prefactors, we arrive to the identity of formal power series in the variables $\mathbf{t}$:
\begin{gather*}
Z_N(\mathbf{t}) = \frac{(-1)^{\frac{N(N - 1)}{2}} \sqrt{\Delta(\boldsymbol{\lambda},\boldsymbol{\lambda})}}{2^N \pi^{\frac{N}{2}} \Delta(\boldsymbol{\lambda})} \operatorname{Pf}_{1 \leq m,n \leq N} \bigl(L_{m,n}\bigr), \\
L_{m,n} = \biggl(\fint_{\mathbb{R}^2} \frac{x - y}{x + y} {\rm e}^{-\frac{1}{2}\lambda_m x^2 - \frac{1}{2}\lambda_n y^2 + V_{\mathbf{t}}(x) + V_{\mathbf{t}}(y)} \,\dd x \dd y\biggr).%\label{ZNL}
\end{gather*}


We would like to rewrite this formula by absorbing the denominator $\Delta(\boldsymbol{\lambda})$ in the Pfaffian. Recall the transposed Vandermonde matrix $W(\boldsymbol{\lambda})^{\mathsf T}$, whose entries are $W(\boldsymbol{\lambda})_{i,n}^{\mathsf T} = \lambda_i^{n}$ indexed by $i \in [N]$ and $n \in \{0,\dots,N - 1\}$. We have
\[
\frac{{\rm Pf}(L)}{\Delta(\boldsymbol{\lambda})} = \frac{{\rm Pf}(L)}{\det W(\boldsymbol{\lambda})^{\mathsf T}} = {\rm Pf}\bigl(\bigl(W(\boldsymbol{\lambda})^{\mathsf T}\bigr)^{-1} LW(\boldsymbol{\lambda})^{-1}\bigr),
\]
and by Cramer's formula for the inverse
\[
\bigl(\bigl(W(\boldsymbol{\lambda})^{\mathsf T}\bigr)^{-1} L W(\boldsymbol{\lambda})^{\mathsf T}\bigr)_{m,n} = v_mv_n\fint_{\mathbb{R}^2} \frac{x - y}{x + y} F_{N;m}(x)F_{N;n}(y) {\rm e}^{V_{\mathbf{t}}(x) + V_{\mathbf{t}}(y)}\,\dd x\dd y,
\]
where $v_n$ are non-zero constants to be chosen later, rows and columns are indexed by $m,n \in \{0,\dots,N - 1\}$, and we introduced
\begin{align*}
F_{N;m}(x) &{}= \frac{1}{v_m}\sum_{i = 1}^{N} \bigl(W(\boldsymbol{\lambda})^{\mathsf T}\bigr)^{-1}_{i,m} {\rm e}^{-\frac{1}{2}\lambda_i x^2} \\
&{}= \frac{\det\bigl(\lambda_i^0 \big| \lambda_i^1  \big| \cdots \big| \lambda_i^{m - 1} \big| {\rm e}^{-\frac{1}{2}\lambda_i x^2} \big| \lambda_i^{m + 1} \big| \cdots \big| \lambda_i^{N - 1}\bigr)}{v_m \Delta(\boldsymbol{\lambda})}.
\end{align*}
With Taylor formula in integral form at order $N$ near $0$, we can write
\begin{gather*}
{\rm e}^{-\frac{1}{2}\lambda_ix^2} = P_{N - 1}\bigl(-\tfrac{1}{2}\lambda_ix^2\bigr) + \frac{\bigl(-\frac{1}{2}\lambda_i x^2\bigr)^m}{m!} + R_{N}\bigl(-\tfrac{1}{2}\lambda_i x^2\bigr), \\
R_N(\xi) = \frac{\xi^{N}}{(N - 1)!} \int_{0}^{1}(1-u)^{N - 1} {\rm e}^{\xi u}\, \dd u
\end{gather*}
for some polynomial $P_{N - 1}$ of degree at most $N - 1$ and without its term of degree $m$ (which we wrote separately). The contribution of $P_{N - 1}$ disappears as it is a linear combination of the other columns, while the contribution of the degree $m$ term simply retrieves the Vandermonde determinant. Hence,
\[
F_{N;m}(x) = \frac{(-1)^mx^{2m}}{2^m m! v_m} + \frac{\det\bigl(\lambda_i^0  \big| \lambda_i^1  \big| \cdots \big| \lambda_i^{m - 1} \big| R_N\bigl(-\tfrac{1}{2}\lambda_ix^2\bigr) \big| \lambda_i^{m + 1} \big| \cdots \big| \lambda_i^{N - 1}\bigr)}{v_m \Delta(\boldsymbol{\lambda})}.
\]
We now choose $v_m = \frac{(-1)^m}{2^m m!}$ to get $F_{N;m}(x) = x^{2m} + O\bigl(x^{2N}\bigr)$ when $x \rightarrow 0$. Introducing the matrix
\[
K_{N;m,n}(\mathbf{t}) = \fint_{\mathbb{R}^2} \frac{x - y}{x + y}F_{N;m}(x)F_{N;n}(y) {\rm e}^{V_{\mathbf{t}}(x) + V_{\mathbf{t}}(y)}\, \dd x \dd y,
\]
we arrive to
\begin{equation*}
\begin{split}
Z_N(\mathbf{t}) & = \frac{(-1)^{\frac{N(N - 1)}{2}} \sqrt{\Delta(\boldsymbol{\lambda},\boldsymbol{\lambda})} \prod_{n = 0}^{N - 1} v_n}{2^N \pi^{\frac{N}{2}}}   \operatorname{Pf}_{0 \leq m,n \leq N - 1} K_{N;m,n}(\mathbf{t}) \\
& = \frac{\sqrt{\Delta(\boldsymbol{\lambda},\boldsymbol{\lambda})}}{2^{\frac{N^2}{2}} (2\pi)^{\frac{N}{2}} \prod_{n = 1}^{N - 1} n!}  \operatorname{Pf}_{0 \leq m,n \leq N - 1} K_{N;m,n}(\mathbf{t}).
\end{split}
\end{equation*}


\begin{thebibliography}{99}
\bibitem[6]{Date:1981qz}
Date E., Jimbo M., Kashiwara M., Miwa T., Transformation groups for soliton
 equations.~{IV}. {A}~new hierarchy of soliton equations of {KP}-type,
 \href{https://doi.org/10.1016/0167-2789(82)90041-0}{\textit{Phys.~D}} \textbf{4} (1982), 343--365.

\bibitem[7]{deBruijn1955}
de~Bruijn N.G., On some multiple integrals involving determinants,
 \textit{J.~Indian Math. Soc. (N.S.)} \textbf{19} (1955), 133--151.

\bibitem[12]{AnaN}
Isaacson E., Keller H.B., Analysis of numerical methods, Dover Publications,
 New York, 1994.

\bibitem[13]{Itzykson:1979fi}
Itzykson C., Zuber J.B., The planar approximation.~{II}, \href{https://doi.org/10.1063/1.524438}{\textit{J.~Math.
 Phys.}} \textbf{21} (1980), 411--421.

\bibitem[20]{Schur1911}
Schur J., \"Uber die {D}arstellung der symmetrischen und der alternierenden
 {G}ruppe durch gebrochene lineare {S}ubstitutionen, \href{https://doi.org/10.1515/crll.1911.139.155}{\textit{J.~Reine Angew.
 Math.}} \textbf{139} (1911), 155--250.

\end{thebibliography}
\end{document}
